% 
% EMSEC seminar template, 24.09.2010
% david.oswald@rub.de
% 
% based on
% EMSEC thesis template, 13.09.2010
% 
% Calls all chapters
% Each chapter is contained in a seperate "\input" file 
% Style: doublesided, scrreport, dina4
%
% Benedikt Driessen, 2010
% benedikt.driessen@rub.de
%

% *************************************************************
% ENTER INFORMATION ABOUT AUTHOR, TITLE and TYPE HERE
\newcommand{\thauthor}{Leon von Mulert}
\newcommand{\thtitle}{Glitches in the Matrix and where to find them}
% *************************************************************

\documentclass[a4paper,12pt,twoside,openany,headsepline,bibliography=totocnumbered]{scrbook}

% Import settings, packages, etc.
%
% List of packages, ...
%

\usepackage{latexsym}
\usepackage{fullpage}
\usepackage{graphicx}
\usepackage{amsmath}
\usepackage{amssymb}
\usepackage{a4}
\usepackage{amsfonts}
\usepackage{eucal}
\usepackage{fancyhdr}
\usepackage[T1]{url}
\usepackage{listings}
\usepackage[printonlyused]{acronym}
\usepackage{algorithmic}
\usepackage{algorithm}
\usepackage{multicol}
\usepackage{hyphenat}
\usepackage{todonotes}

% Urlsyle
\urlstyle{tt}
% Parskip
\setlength{\parskip}{1ex plus0.5ex minus0.2ex}
% Seperation between list items
\setlength{\itemsep}{0ex plus0.2ex}

% Fancy headers
\setlength{\headsep}{8mm}
\pagestyle{fancyplain}
\renewcommand{\chaptermark}[1]{\markboth{#1}{#1}}
\renewcommand{\sectionmark}[1]{\markright{\thesection\ #1}}
\lhead[\fancyplain{}{\thepage}]{\fancyplain{}{\slshape \rightmark}}
\rhead[\fancyplain{}{\slshape \leftmark}]{\fancyplain{}{\thepage}}
\cfoot{}



\begin{document}

% Import frontpage
%
% Frontpage ripped from S. Rinne
%

\frontmatter

\begin{titlepage}
 \enlargethispage{3cm}
 \vspace*{-32mm}\hspace*{120mm}
 \includegraphics[scale=1.0]{rub_logo.eps}
 
 \vspace*{11cm}\hspace*{0mm}
 \begin{minipage}[b]{1\linewidth}
  \sffamily
  \hspace{-17.2mm}\includegraphics[scale=1.0]{rub_slogan.eps}\\
  
  \nohyphens{
   {\bfseries \LARGE \sffamily {\thtitle}}
  }\\
  
  \large{
   \thauthor
  }\\
  
  \vspace*{35mm}
  \normalsize{
   Seminarausarbeitung\\
   \today\\
   Chair for Security Engineering - Prof. Dr.-Ing. Tim G{\"u}neysu\\
	 Advisor: Florian Stolz
  }
 \end{minipage}
\end{titlepage}


\newpage\thispagestyle{empty}


% Abstract
\section*{Abstract}
The following text\footnote{Shamelessly ripped from \url{http://www.leeds.ac.uk/educol/abstract.htm}} describes what your abstract should be about.

The abstract should convey to the reader concisely and accurately within the space of a few sentences, the claim to knowledge that the authors are making. It should indicate the boundaries of space and time within which the enquiry has occurred. If there is a claim to generality beyond the boundaries of the enquiry the basis of that claim should be given, for example that a random sample is thought to be representative of a larger population. There should also be a hint of the method of enquiry.

The boundaries of an enquiry are important - and are unfortunately too often omitted from abstracts. This is due to the regrettable tendency for researchers to generalise their results from, for example, a few schools to all schools, and to imply that what was true at a particular time, is true for all time. Some reference to the geographical location of the children, or teachers, or schools on whom the claim to knowledge rests should be made. Because of the international nature of the research community it is worth making clear in what country the research took place. Also the period in which the data was collected should be stated.

The abstract should be a condensation of the substance of the paper, not a trailer, nor an introduction. Journals and thesis regulations usually put a limit of around 200 to 300 words to the length of an abstract. “Trailer” is a term borrowed from the cinema industry to describe a showing of a few highlights in order to win an audience. An “Introduction” tells that something is coming, but doesn’t reveal its substance. These are not what is needed.

Abstracts are recycled in abstract journals and electronic networks and provide the main vehicle for other researchers to become aware of particular studies. Hence the more clearly they convey the claim to knowledge of the original paper the more useful they are in helping the reader to decide whether it is worth taking the trouble to obtain and read the original and possibly cite it in his/her own writing.

Both the abstract and the paper should make sense without the other.
\clearpage

\tableofcontents
\mainmatter

% List of acronyms
\chapter*{Acronyms}
\begin{acronym}
 \setlength{\itemsep}{0.2em}
 \acro{ACR}{ACRONYM1}
 % etc.
\end{acronym}

%
% Include all your chapters as .tex files, each file contains sections \section{name of section},
% subsections \subsections{...} and so on...
%

\pagenumbering{arabic} 

%
% Sample introduction of your thesis
%

\chapter{Introduction}
Physical access to otherwise secure systems has become much more common due to the widespread use of embedded devices in today's world. This circumstance enables attacks that are impossible in more restricted scenarios and thus, the need for countermeasures against such physical attacks. 

Indeed, glitching a system, i.e., introducing a physical disturbance such that the execution of the processor is influenced, is one of the many ways attackers may be able to leverage this increased level of access. These hardware-induced faults have already been used to compromise commercially available systems, most notably the XBOX 360\,\cite{free60} and PlayStation 3\,\cite{ps3hyper}. Similar to Side-Channel Analysis\,\cite{standaert2010introduction}, glitching-based attacks offer a new angle of attack against systems that are running semantically secure software without any conventional vulnerabilities.

Although fault attacks against cryptographic computations have been studied extensively in the past\,\cite{boneh1997importance}\todo{cite more}, precise glitching enables much more devious attacks, such as gaining \texttt{root} access by flipping only a single bit in the \texttt{sudo} binary\,\cite{gruss2018another}. Instead of merely faulting the computation of an important value, glitches are used to alter the executed instructions, e.g., altering the condition on a branch instruction, and hence violate the control flow integrity of the attacked program.  As a result of this, many of the suggested defences against fault attacks, e.g., redundant computations\cite{chang2019evaluating}, can be skipped or circumvented entirely. After all, why forge a signature when you can simply skip the signature check?\todo{too informal}

In this paper we introduce a formal description of hardware-induced faults and give an overview of \todo{Von anderen Papern vorgeschlagen??} defences suggested by related works. We categorise these countermeasures by whether they are software-only, hardware-only or require a hybrid solution and evaluate them within those categories, based on their quality (error detection rate and latency) and their induced overhead in memory or runtime.

\todo[inline]{Short overview of attack categories (control flow violation vs. faulty computation)}
\todo[inline]{Short overview of defence categories (control flow integrity via signatures; redundant computations)}
%\input{theory}
%\input{design}
%\input{implementation}
%\input{results}
%
% Sample conclusion of your thesis
%

\chapter{Conclusion}
Conclude your thesis and discuss your results...


% Generate list of figures
\newpage
\listoffigures

% Generate list of tables
\newpage
\listoftables
\clearpage

% Generate bibliography with bibtex, the bibfile here is "bibliography.bib", use alphanumerical style
\begin{appendix}
 \bibliography{bibliography}
 \bibliographystyle{alpha}
\end{appendix}

\end{document} 
